% Lojban Dictionary
% Created by Leo Molas
% leos(dot)molas(at)gmail(dot)com
% http://github.com/leosmolas

% This work is licensed under the Creative Commons 
% Attribution-Noncommercial-Share Alike 3.0 Unported License. To view a 
% copy of this license, visit 
% http://creativecommons.org/licenses/by-nc-sa/3.0/ or send a letter to 
% Creative Commons, 171 Second Street, Suite 300, San Francisco, 
% California, 94105, USA. 

% You can find the latest version of this .tex, and the others that are
% used, at http://github.com/leosmolas/pyvlacku/tree/master/pysrc/

\documentclass[ipa,twoside]{report} 

\usepackage[english]{babel}
\usepackage[phonetic]{lexikon}
\usepackage[latin1]{inputenc}
\usepackage{fullpage}
\usepackage{multicol}
\usepackage[colorlinks=true]{hyperref}
\usepackage{fancyhdr}
\usepackage{hyphenat}
\usepackage{graphicx}
\usepackage{array}
\usepackage{eurosym}

\selectlanguage{english}

\author{The Lojban Community}
\title{Lojban - English\\English - Lojban}

\begin{document}

%this ones are the fields of the dictionary, defined in lexikon.
\setlength{\leftfield}{0.15\textwidth}
\setlength{\rightfield}{0.3\textwidth}

%I redifine plain style to make it like fancy, but without the header; just the page number
\fancypagestyle{plain}{%
	\fancyhf{}%
	\fancyfoot[LE,RO]{\thepage}%
	\renewcommand{\headrulewidth}{0pt}
	\renewcommand{\footrulewidth}{0pt}
}

\maketitle

\tableofcontents


\chapter*{Version history}
\addcontentsline{toc}{chapter}{Version history}
\label{cha:history}
\paragraph*{Version v1.0}

The first version of the dictionary, fully functional. If you find anything wrong, let me know.

%I don't have the skills to write this well. If you can, do it :) and let me know

\chapter*{How to use this dictionary}
\addcontentsline{toc}{chapter}{How to use this dictionary}
\label{cha:howto}
In the chapter \textbf{Lojban - English}, you will find the list of all the \textsl{valsi}, in alphabetical order. It will be look like:

\begin{center}
\begin{minipage}[t]{0.55\textwidth}
	\dictchar{cy.}
	\dictentry{ce}{/she/}{c}{\textit{cec}\\\textbf{JOI}} %cambiar!!!
 	{}{}{}{non-logical connective: set link, unordered; "and also", but forming a set.}
\end{minipage}
\end{center}

In the left column, you can find the \textsl{valsi} that is being defined. It will be followed by a letter, that stands for the type of the word, as in this list:

\begin{tabular}{ l l }
  \textit{g} & \textbf{\emph{gismu}}  \\
  \textit{f} & \textbf{\emph{fu'ivla}}  \\
  \textit{l} & \textbf{\emph{lujvo}}  \\
  \textit{c} & \textbf{\emph{cmavo}} \\
  \textit{n} & \textbf{\emph{cmene}} \\
  \textit{cc} & \textbf{\emph{cmavo cluster}} \\
  \textit{eg} & \textbf{\emph{experimental gismu}} \\
\end{tabular}

Below the \textsl{valsi}, in between slashes (``/'') you will find the word pronunciation, in the \textbf{IPA} representation. After it, in \textit{italics}, can appear the \textsl{rafsi} associated with the \textsl{valsi}. If the word is a \textsl{cmavo} or a \textsl{cmavo cluster}, it can appear in \textbf{bold}, the \textsl{selma'o} that the \textsl{valsi} belongs to.

In the right column, you will find the definition in English, with its \textbf{place structure}. Below it, in \textit{italics}, can appear a comment about the word.
%
%\addcontentsline{toc}{chapter}{}
%\chapter*{Version history}
%\label{cha:history}

%I don't know if this is working. I should change it...
\setlength{\headheight}{25pt}
% here, I redifine the fancyhdr like a dictionary, with the first and last words in the header. The package "lexikon" does the magic here, assigning the words
\fancyhead{}					 		% empty out the header
\fancyfoot{}							% empty out the footer
\fancyhead[LE,LO]{\rightmark}			% left side, odd and even pages
\fancyhead[RE,RO]{\leftmark}			% right side, odd and even pages
\fancyfoot[LE,RO]{\thepage}				% left side even, right side odd

%this is the line between the columns
\setlength{\columnseprule}{0.3pt}
\pagestyle{fancy}
\thispagestyle{plain}

\chapter*{Lojban - English}
\addcontentsline{toc}{chapter}{Lojban - English}

\label{cha:lojen}
\begin{multicols}{2}
%this is the file generated with the python script
\input{lojen}
  
\end{multicols}

\newpage 
\setlength{\leftfield}{0.22\textwidth}
\setlength{\rightfield}{0.22\textwidth}

\thispagestyle{plain}

%the inverse dictionary
\chapter*{English - Lojban}
\addcontentsline{toc}{chapter}{English - Lojban}

\label{cha:enloj}
\begin{multicols}{2}
% this is the other file generated with the python script
\input{enloj}

\end{multicols}

% boring legalities
\chapter*{Copyleft}
\addcontentsline{toc}{chapter}{Copyleft}

\pagestyle{plain}

\section*{Lojban}
\addcontentsline{toc}{section}{Lojban}

Lojban words and grammar are in \textbf{Public Domain}. You can find the latest in the \textbf{Complete Lojban Language}%
\footnote{You can download it or purchase it, in \url{http://www.lojban.org/tiki/The+Complete+Lojban+Language}}.

All the words were extracted from the \textsl{jbovlaste} web(\url{http://jbovlaste.lojban.org/}), under the following licence\footnote{Extracted from \url{http://www.lojban.org/tiki/LLG+Web+Copyright+License&bl}}.

\begin{quotation}
CONTACT THAT ORGANIZATION AT: 2904 BEAU LANE FAIRFAX VA 22031 USA 703-385-0273 lojban@lojban.org

PERMISSION TO COPY INDIVIDUAL DOCUMENTS GRANTED SUBJECT TO YOUR VERIFICATION THAT YOU HAVE THE LATEST VERSION OF THE DOCUMENT IN QUESTION, THAT YOUR DISTRIBUTION BE FOR PROMOTION OF LOJBAN, THAT THERE IS NO CHARGE FOR THE PRODUCT, AND THAT THE COPYRIGHT NOTICE IS INCLUDED INTACT IN THE COPY.

PERMISSION GRANTED FOR USE IN DERIVED WORKS PROVIDED THAT THE FACTS OF DERIVATION IS STATED, AND PROVIDING THE NAME AND ADDRESS OF THE LOGICAL LANGUAGE GROUP, INC. IS PROVIDED AS A SOURCE OF FURTHER INFORMATION ABOUT LOJBAN.

In accordance with the Logical Language Group, Inc. policy, certain Lojban language definition materials are placed irrevocably in the Public Domain.

Individual documents may contain more precise copyright statements.

signed:

Robert LeChevalier President The Logical Language Group, Inc.

This material is being made available without charge by the Logical Language Group for the purpose of broadening interest in the use of the constructed human language ``Lojban - A Realization Of Loglan''.

You are, however, urged to send contributions to the above address in support of the development of the Lojban language. We are a 501(c)(3) non-profit organization and such donations are tax-deductible in the US, to the extent they exceed the price of things we send you. While it costs us little for you to get a copy of the files on the archives, the network is not free. la lojbangirz. spent some \$1500 during 1999 on network access, supporting Lojban List, the archives, and responding to people like you. We do not make a profit on our printed materials, so it is your donations that will allow us to continue to serve the network community of Lojbanists and people curious about Lojban.
\end{quotation}

\section*{This book}
\addcontentsline{toc}{section}{This book}

\begin{center}
	\includegraphics{img/by-nc-sa.png}
\end{center}

\begin{quotation}
This work is licensed under the Creative Commons Attribution-Noncommercial-Share Alike 3.0 Unported License. To view a copy of this license, visit \url{http://creativecommons.org/licenses/by-nc-sa/3.0/} or send a letter to Creative Commons, 171 Second Street, Suite 300, San Francisco, California, 94105, USA.
\end{quotation}

So...
\begin{quotation}

\subsection*{You are free}

% this two definitions here are... well, I don't understand them. I extracted them from 
% http://ask.metafilter.com/60851/How-to-center-vertically-in-a-LaTeX-table
\newcolumntype{S}{>{\arraybackslash} m{.8\linewidth} }
\newcolumntype{T}{>{\arraybackslash} m{4em }}
\begin{tabular}{T S}%{c c}
  \includegraphics[width=3em]{img/sharelarge.png} & \textbf{to Share} --- to copy, distribute and transmit the work  \\
  \includegraphics[width=3em]{img/remixlarge.png} & \textbf{to Remix} --- to adapt the work  \\
\end{tabular}

\subsection*{Under the following conditions}

\begin{tabular}{ T S }
  \includegraphics[width=3em]{img/bylarge.png} & \textbf{Attribution} --- You must attribute the work in the manner specified by the author or licensor (but not in any way that suggests that they endorse you or your use of the work). \\
  \includegraphics[width=3em]{img/nclarge.png} & \textbf{Noncommertial} --- You may not use this work for commercial purposes.  \\
  \includegraphics[width=3em]{img/sharelarge.png} & \textbf{Share Alike} --- If you alter, transform, or build upon this work, you may distribute the resulting work only under the same or similar license to this one. \\
\end{tabular}

\end{quotation}

The \textbf{Noncommertial} clause has one exception: The \textbf{Logical Language Group} (a.k.a. LLG) have my total permission to use sell this book, in any form.

You can change the book, as long as it remains in an open licence, but it's going to be difficult, since it's a .pdf. You can find the source code of this book(\texttt{.tex}) at \url{http://github.com/leosmolas/pyvlacku/tree/master/pysrc/}. 

But you might want to know how I created this. That lead us to the next section. 

\section*{The Software (pyvlacku)}
\addcontentsline{toc}{section}{The Software (pyvlacku)}
This book was created with the help of a software developed by \textbf{Leo Molas} (or just \textsl{.leos.}), called \textbf{pyvlacku}, and it's under GPL v3.0. You can find the source code at \url{http://github.com/leosmolas/pyvlacku}. 

\end{document}
